\section{Implementation}
\label{sec:implementation}

All implementation work described in this section is the author's. It lives on the branch
\path{feat/phase1-notification-coalescing} of the project repository and is submitted as
pull request~\#1~\cite{tap-pr1}. Relative to the branch point (\code{e52bb66}, the
upstream \code{0.2.0-beta.7} release), the implementation commits touch 91~files with
approximately 8{,}200 lines added and 140 removed, of which roughly 3{,}200 lines are
production source, 4{,}300 lines are tests, and the remainder is documentation.

\subsection{Increments}

The work was landed as five increments, each of which left the repository lint-clean,
type-clean and green under \code{bun run test}, in the dependency order that principle~P2
prescribes. Table~\ref{tab:increments} summarises them; the full commit list is in
Appendix~\ref{app:commits}.

\begin{table}[htbp]
\centering\small
\caption{Implementation increments, in landing order.}
\label{tab:increments}
\begin{tabularx}{\textwidth}{@{}l>{\raggedright\arraybackslash}X>{\raggedright\arraybackslash}X@{}}
\toprule
\textbf{Increment} & \textbf{Scope} & \textbf{Principal modules} \\
\midrule
1. Coalescing &
Coalesce at enqueue; escalation-first render plan with grouped overflow summary; flood demo script; skill documentation of the new block semantics. &
\path{tapd/notification-queue.ts}, \path{tapd/notification-format.ts}, \path{scripts/demo-notification-flood.ts} \\
2. Ledger &
File-backed per-peer attention ledger; accounting on drain derived from the canonical render plan; \code{tap attention show}. &
\path{core/attention/ledger.ts}, \code{tapd/daemon.ts}, \path{cli/commands/attention-show.ts} \\
3. Pricing &
\code{attention} block in config and registration file (validated on both sides); price discovery in \code{identity resolve}; \code{--dry-run} cost preview; \code{-32050} rejection of un-granted \code{message/send}. &
\code{core/config/*}, \path{core/identity/registration-file.ts}, \path{core/common/errors.ts}, \path{core/runtime/service.ts}, \path{core/transport/xmtp.ts} \\
4. Postage &
Credit ledger (issued/held), stamps, credit certificates; \code{postage/topup} and \code{postage/balance} handlers; auto-stamping in \code{sendMessage}; top-up orchestration; \code{tap postage topup/balance}; \code{@trustedagents/app-postage}; \code{tapd} routes. &
\code{core/postage/*}, \code{app-postage/}, \path{cli/commands/postage-*.ts}, \path{tapd/http/routes/postage.ts} \\
5. Priority \& quota &
Priority tier on stamps; escalation classification with 400-char excerpt; wake-up in OpenClaw; \code{notificationsPerWeek} folding at drain; \code{--priority} threaded through CLI, \code{tapd}, OpenClaw tool and Hermes plugin. &
\path{tapd/notification-quota.ts}, \path{tapd/notification-classifier.ts}, \path{openclaw-plugin/escalation-watcher.ts}, \path{cli/assets/hermes/plugin/} \\
\bottomrule
\end{tabularx}
\end{table}

\subsection{Data structures}

\paragraph{Notifications.}
\code{TapNotification} gained three fields (\code{coalesceKey}, \code{coalesceStrategy}
and \code{count}), all optional so that a batch produced by an older daemon renders
unchanged. The queue keeps a \code{Map<coalesceKey, entry>} beside
the array buffer; the invariant that the map never points at an entry outside the buffer
is maintained at exactly two places (drain and max-size eviction) and is covered by unit
tests.

\paragraph{Attention ledger file.}
Listing~\ref{lst:ledger} shows the on-disk shape. Day buckets are keyed by UTC date so
that the weekly window (\code{renderedInWindow(peer, 7)}) is a simple prefix comparison
over keys, and pruning on write bounds the file at the retention period.

\begin{lstlisting}[language=json,caption={Shape of \code{attention-ledger.json} (abridged).},label={lst:ledger}]
{
  "version": 1,
  "identity": { "chain": "eip155:8453", "agentId": 1 },
  "days": {
    "2026-08-27": {
      "peers": {
        "eip155:8453#11": { "notificationsRendered": 1, "tokensInjected": 24,
                            "escalations": 0, "overflowSuppressed": 0 },
        "unattributed":   { "notificationsRendered": 1, "tokensInjected": 31,
                            "escalations": 1, "overflowSuppressed": 0 }
      },
      "totals": { "notificationsRendered": 2, "tokensInjected": 55,
                  "escalations": 1, "overflowSuppressed": 0 }
    }
  }
}
\end{lstlisting}

\paragraph{Postage state.}
\path{<dataDir>/apps/postage/state.json} holds two maps keyed by \code{creditId}:
\code{issued} (credits peers bought \emph{with} this agent; fields \code{peer},
\code{amount}, \code{spent}, \code{lastSeq}, \code{lastStampKey}, \code{txHash},
\code{certificate}) and \code{held} (credits this agent bought \emph{at} peers; fields
\code{nextSeq}, \code{certificateVerified} in place of \code{lastSeq}). Amounts are
decimal strings with at most six places and are converted to integer micro-units
(\code{bigint}) for every arithmetic operation; no floating point touches money.

\subsection{Wire formats}

\paragraph{Registration file.}
The \code{trustedAgentProtocol.attention} block (Listing~\ref{lst:regfile}) is validated
by the same \code{validateRegistrationFile} that already checks the XMTP endpoint, on both
the publishing and the resolving side, so a malformed price list is rejected before it
can be cached.

\begin{lstlisting}[language=json,caption={Attention block in the ERC-8004 registration file.},label={lst:regfile}]
"trustedAgentProtocol": {
  "agentAddress": "0x...",
  "attention": {
    "version": "1.0",
    "currency": "USDC",
    "chain": "eip155:167000",
    "pricing": { "grantHolder": "0", "standard": "0.001", "priority": "0.01" }
  }
}
\end{lstlisting}

\paragraph{Rejection.}
The transport previously collapsed any service error into the generic JSON-RPC
\code{-32603}. It now preserves \code{AttentionPaymentRequiredError} as \code{-32050}
with structured \code{data} (Listing~\ref{lst:reject}), and the sending side raises a
\code{TransportRpcError} that keeps \code{rpcCode} and \code{rpcData} so that callers
can react programmatically. The CLI and \code{tapd} translate the same object to HTTP
\code{402 attention\_payment\_required}.

\begin{lstlisting}[language=json,caption={A \code{-32050} rejection as seen by the sender.},label={lst:reject}]
{
  "jsonrpc": "2.0", "id": "req-...",
  "error": {
    "code": -32050,
    "message": "attention payment required: ...",
    "data": {
      "attention": { "version": "1.0", "currency": "USDC",
                     "chain": "eip155:167000",
                     "pricing": { "standard": "0.001", "priority": "0.01" } },
      "postage": { "reason": "insufficient_credit", "creditId": "...",
                   "remaining": "0", "required": "0.001" }
    }
  }
}
\end{lstlisting}

\paragraph{Stamp.}
A stamp is a \code{postage} object with three fields, \code{creditId}, \code{seq} and
\code{cost}, inside the existing \code{trustedAgent} metadata of a \code{message/send}
(for example \code{\{"creditId": "c-7f3a", "seq": 3, "cost": "0.001"\}}).
\path{extractPostageStamp} guards only shape; economic validity is the ledger's decision,
so a malformed stamp can never crash enforcement.

\subsection{The service gate}

The gate is a single block in \code{TapMessagingService.processInboundRequest}, placed
after contact resolution and duplicate detection but before the conversation log. Its
ordering constraints are the subtle part of the implementation and are worth stating:

\begin{enumerate}
  \item Enforcement runs only for \code{message/send} and only when
        \code{config.attention.enforce} is true.
  \item A rejected request marks its journal entry \code{completed} \emph{before} throwing,
        so that a full-history sync hits the duplicate short-circuit rather than re-running
        enforcement per replay.
  \item An accepted stamp completes the journal entry \emph{after} the conversation log.
        A crash between the debit and that completion redelivers the message with its
        journal entry pending; the ledger's \code{lastStampKey} lets that exact redelivery
        pass idempotently, and the event is emitted even though the delivery is a duplicate,
        because otherwise the receiver would keep the money and drop what it bought.
  \item For a grant holder the gate does not reject; it calls
        \code{chargeOptionalPriorityStamp}, which debits only if the stamp meets the
        \code{priority} price and swallows every failure.
\end{enumerate}

The \code{message.received} event carries an optional \code{postage: \{tier, cost\}}
reference only when a stamp was actually debited, so downstream consumers cannot
mistake a free delivery for a paid one.

\subsection{Top-up orchestration}

\code{topUpPostage} runs under the service's execution mutex and inside a transport
session. It (i)~validates the amount; (ii)~pays through the host's
\code{executeTransfer} hook---the same hook \code{tap transfer} uses, so the payment is an
ordinary USDC transfer on the contact's chain, signed through OWS; (iii)~records the held
credit; (iv)~sends the \code{postage/topup} action and waits for the result with a waiter
that settles on \emph{any} result for the request, not only one carrying the
\code{actionId}, so that a pre-postage peer's \code{UNSUPPORTED\_ACTION} reply is read as
a definitive rejection rather than a timeout; (v)~verifies the returned certificate
against the contact's agent address and marks the held credit verified. A transport
failure after payment is returned as \code{status: pending} with the \code{txHash}
rather than thrown, because the payment has happened and the caller must be told so.

\subsection{Host surfaces}

Following the ``thin plugin, fat CLI'' rule, only operations that need the live transport
were added to the plugin surface. Table~\ref{tab:surfaces} lists the new commands.

\begin{table}[htbp]
\centering\small
\caption{Operator-facing surfaces added by the project.}
\label{tab:surfaces}
\begin{tabularx}{\textwidth}{@{}l>{\raggedright\arraybackslash}Xl@{}}
\toprule
\textbf{Surface} & \textbf{Behaviour} & \textbf{Transport?} \\
\midrule
\code{tap attention show} & Per-peer attention usage over the retained window, highest token spend first; reads \code{attention-ledger.json}. & No \\
\code{tap identity resolve <id>} & Now prints the peer's advertised \code{attention} block. & No \\
\code{tap register update} & Publishes \code{attention.pricing} from config into the registration file. & No (on-chain) \\
\code{tap message send --dry-run} & Shows the peer's price, \code{postage\_remaining} and \code{would\_stamp}; sends nothing. & No \\
\code{tap message send --priority} & Pays the peer's \code{priority} price so the message escalates and may wake the peer. & Yes \\
\code{tap postage topup <peer> --amount} & Pays USDC on-chain, opens a credit, stores the certificate; \code{--yes}, \code{--dry-run}. & Yes \\
\code{tap postage balance [--peer]} & Credits held at peers and issued to peers. & No \\
\code{tap\_gateway send\_message \{priority\}} & OpenClaw and Hermes tool flag mirroring \code{--priority}. & Yes \\
\code{tapd} \code{POST /api/postage/*}, \code{/api/notifications/drain} & Daemon routes behind which the CLI and plugins operate. & --- \\
\bottomrule
\end{tabularx}
\end{table}

The Hermes Python plugin, which mirrors the render plan rather than importing it, was
extended to render the \code{(xN)} suffix and to forward the \code{priority} flag, and its
Python tests (\code{test\_client.py}) pin both behaviours; the folded-quota summary needs no
host-side change because it arrives as an ordinary \code{summary} notification.

\subsection{Testing}

Testing follows the repository's existing three tiers.

\paragraph{Unit.} New suites cover the queue invariants (merge in place, index pruned on
eviction, \code{replace} strategy never counts), the render plan (partition stability,
grouped summary text, cap), the attention ledger (bucketing, retention, window counts),
the postage ledger (idempotent \code{recordIssued}, one credit per \code{txHash}, per-peer
cap, replay rejection, redelivery exception, no double spend under concurrent debits), the
certificate (sign/verify round trip; wrong signer rejected), the payload parsers, the
quota fold (escalations excluded; fail-open on lookup error), and the OpenClaw escalation
watcher.

\paragraph{Service.} \code{service.test.ts} exercises the gate against a mocked transport:
enforcement off is a no-op; un-granted and unstamped is rejected with a quote; a grant
exempts; a valid stamp debits; an exhausted credit yields \code{insufficient\_credit}; a
grant holder's priority stamp buys the escalation tier.

\paragraph{Mock end-to-end.} Three new phases were added to the canonical scenario list
shared by both E2E suites (Table~\ref{tab:e2e}). They run over the loopback transport on
every pull request. They are deliberately \emph{not} mirrored into the live-mainnet suite
yet: the key assertions are negative (``B never logged the message''), which are
timing-flaky over real XMTP, and both mechanisms are opt-in so live behaviour is
unchanged. This gap is recorded as a follow-up in the pull request.

\begin{table}[htbp]
\centering\small
\caption{Mock E2E scenarios added (phases 6--8 of \code{packages/cli/test/e2e/scenarios.ts}).}
\label{tab:e2e}
\begin{tabularx}{\textwidth}{@{}l>{\raggedright\arraybackslash}X@{}}
\toprule
\textbf{Phase} & \textbf{Scenarios} \\
\midrule
6 Attention &
Preview cost with \code{--dry-run}; enable enforcement on B; send without grant rejected with quote; \code{message/send} grant exempts sender. \\
7 Postage &
Revoke grant so postage applies; top-up buys credit with signed certificate; auto-stamped message delivers and debits; exhausted credit rejected with top-up quote; \code{postage balance} shows held and issued credits on both sides. \\
8 Wake \& quota &
Publish a \code{priority} price; priority stamp escalates with a 400-char excerpt; over-quota grant holder's chatter folds to one summary line; grant holder's priority stamp still buys the wake-up. \\
\bottomrule
\end{tabularx}
\end{table}

\subsection{Documentation}

Every new command and flag was added to the canonical skill file
\path{skills/trusted-agents/SKILL.md}, which the OpenClaw plugin and the Hermes installer
copy at build time, and the \code{notificationsPerWeek} constraint to
\path{references/permissions-v1.md}. The skill's ``how to read the block'' guidance was
rewritten so that an agent encountering \code{(xN)}, a folded-quota summary or a
``Priority message from'' escalation knows what each means and how to act. The
repository's \code{README} and \code{Agents.md} were aligned in a final documentation
commit.
